\documentclass[11pt]{article}
\usepackage[T1]{fontenc}
\usepackage{lmodern}
\usepackage[margin=1in]{geometry}
\usepackage{amsmath,amssymb,amsthm,mathrsfs,booktabs,graphicx,microtype}
\usepackage[colorlinks=true,linkcolor=blue,urlcolor=blue,citecolor=blue]{hyperref}
\newtheorem{proposition}{Proposition}
\newtheorem{corollary}{Corollary}
\newcommand{\eps}{\varepsilon}
\newcommand{\avg}[1]{\left\langle #1\right\rangle}
\newcommand{\Q}{\mathbb Q}
\title{Quadratic Einstein--scalar sources and an intrinsic material-shell diagnostic}
\author{Ricardo Maldonado\\Independent Researcher\\Henderson, Nevada, United States}
\date{Reproducible working note, 16 October 2026 UTC}
\begin{document}
\maketitle
\begin{abstract}
For a retained conditional Einstein--scalar problem with a timelike material
shell, we complete a component evaluator for the quadratic Ricci, natural
Einstein and relative-harmonic reduced sources and reconstruct their complete
23-amplitude angular projection. Exact polynomial identities retain both
transverse-traceless angular orientations. Authenticated differentiated
first-field error bounds give ordinary metric and scalar two-jet enclosures,
and outward interval evaluation bounds the actual quadratic forcing on the
full causal collar. We also define an intrinsic shell scalar whose background
and first variations vanish. Its sphere-mean second parameter derivative has
exactly two simple zeros in each of two certified time brackets. A separate
quantitative consequence of the retained second-response derivative bound
proves persistence of two scalar-mean signs at positive times. The universal
curvature and gauge formulas are established results; the contribution here
is their explicit model-specific evaluation and reproducible certification.
No complete numerical second waveform, finite-amplitude nonlinear family,
invariant force or cosmological-origin conclusion is established.
\end{abstract}

\section{Conditional problem and conventions}

The background and the selected first and second initial prescriptions are
those of the preceding public computational checkpoint~\cite{parent}.
They remain premises here. In five spacetime dimensions write
\begin{equation}
 g_0=dy^2+\rho(y)^2[-dt^2+c(t)^2\kappa],\qquad
 c=\cosh t,\quad \theta=\tanh t,\quad a=\rho_y/\rho,\quad
 p=\phi_{0,y},
\end{equation}
where $\kappa$ is the unit three-sphere metric. The potential $\mathcal V$
and scalar-dependent shell law are fixed by that checkpoint. The scaled
collar is
\begin{equation}
 \begin{gathered}
 \ell=10^{-6},\quad x=t/\ell,\quad r=z/\ell,\quad dy=-\rho\,dz,\\
 D_{\rm base}=\{0\le x\le3,\ r\ge0,\ x+r\le8\},\qquad
 D=D_{\rm base}\times S^3.
 \end{gathered}
 \label{eq:domain}
\end{equation}
Every second variation denotes a second parameter derivative:
$g_\eps=g_0+\eps q+\eps^2 k/2+o(\eps^2)$ and
$\phi_\eps=\phi_0+\eps\eta+\eps^2 w/2+o(\eps^2)$.
Thus a displayed quadratic Hessian is twice the coefficient of $\eps^2$.

Use the orthonormal background frame $e_y=\partial_y$,
$e_0=\rho^{-1}\partial_t$, $e_a=(\rho c)^{-1}E_a$, with $e_0$ timelike.
For $Y=n_1^2-n_2^2$ the actual first harmonic representative has sectors
\begin{equation}
 \begin{gathered}
 q_{yy}=AY,\quad q_{y0}=BY,\quad q_{ya}=C D_aY,\quad
 q_{00}=DY,\quad q_{0a}=E D_aY,\\
 q_{ab}=LY\kappa_{ab}+M D_aD_bY,\quad \eta=sY.
 \end{gathered}
 \label{eq:amplitudes}
\end{equation}
Here the letter $D$ in $q_{00}$ is an amplitude, while $D_a$ denotes the
sphere derivative. In terms of the retained master amplitudes $u,v$ and raw
gauge amplitudes $f,b,d$,
\begin{align}
 A&=-4u+2f_y,& B&=f_t/\rho-\rho b_y,\\
 C&=f/(\rho c)+\rho c d_y,& D&=-2(u+af+b_t),\\
 E&=c d_t-b/c,& L&=2(u+af+\theta b),\\
 M&=2d,& s&=v+pf.
\end{align}
The actual first solve supplies $C_1=DC[q]=0$ as a field, including its
derivative consequences. Arbitrary manufactured jets need not obey this
condition. On the material shell $f=0$ and $b_y=d_y=0$; tangential derivatives
of these identities are retained. No new shell displacement is prescribed.

\section{Complete geometric quadratic sources}

The natural residuals are
\begin{equation}
 \mathscr E_{AB}=R_{AB}-\phi_A\phi_B-\tfrac23\mathcal V(\phi)g_{AB},
 \qquad \mathscr S=\Box_g\phi-\mathcal V'(\phi).
\end{equation}
All following contractions and unmarked derivatives use $g_0$. Set
\begin{align}
 P^A{}_{BC}&=\tfrac12g_0^{AD}
 (\nabla_Bq_{CD}+\nabla_Cq_{BD}-\nabla_Dq_{BC}),\qquad
 J^A{}_{BC}=-2q^A{}_D P^D{}_{BC},\\
 Q^{\rm Ric}_{AB}&=\nabla_CJ^C{}_{AB}-\nabla_BJ^C{}_{AC}
 +2(P^C{}_{CD}P^D{}_{AB}-P^C{}_{BD}P^D{}_{AC}),\\
 Q^E_{AB}&=Q^{\rm Ric}_{AB}-2\eta_A\eta_B
 -\tfrac23\big[\mathcal V''\eta^2(g_0)_{AB}
                     +2\mathcal V'\eta q_{AB}\big].
 \label{eq:QE}
\end{align}
These are the standard affine-path curvature variations, in parameter
derivative normalization~\cite{bmm2006,xpert}. The source calculation acts
on every tensor slot, including the derivative slot of $\nabla q$.
Ordinary derivatives of the frame amplitudes are converted before the
connection is applied; they are not themselves covariant jets.

\subsection{Gauge lowering and an off-shell check}

Let $C^A(g)=g^{BC}(\Gamma[g]-\Gamma[g_0])^A{}_{BC}$ and
$C_B(g)=g_{BA}C^A(g)$. Along $g_\eps=g_0+\eps q$, the exact connection
difference gives
\begin{equation}
 C_B(g_\eps)=\eps g_\eps^{CD}P_{BCD},\qquad
 \widehat Q_{C,B}:=\frac{d^2}{d\eps^2}C_B(g_\eps)\big|_0
 =-2q^{CD}P_{BCD}.                                      \label{eq:lower}
\end{equation}
This identity is off shell. By contrast, the upper-vector Hessian is
$Q_C^A=-2q^A{}_D C_1^D-2q^{CD}P^A{}_{CD}$.
The varying-lowering term $2q_{BA}C_1^A$ cancels the first term when forming
\eqref{eq:lower}; background lowering alone does not.

For a reduction using $\nabla^g_{(A}C_{B)}$, the general affine Hessian is
\begin{equation}
 Q^{E,\mathrm{red,gen}}_{AB}=Q^E_{AB}
       -\nabla^0_{(A}\widehat Q_{C,B)}+2P^C{}_{AB}(C_1)_C.
 \label{eq:generalred}
\end{equation}
Parentheses contain one half. The last term checks differentiation of the
covariant derivative. For the registered harmonic field $C_1=0$,
\begin{equation}
 Q^{E,\mathrm{red}}_{AB}=Q^E_{AB}-\nabla^0_{(A}Q_{C,B)},\qquad
 Q_{C,B}=-2q^{CD}P_{BCD}.                              \label{eq:actualred}
\end{equation}
The evaluator exposes both \eqref{eq:generalred} and \eqref{eq:actualred}.
Off-shell test jets are compared using the general formula. Enclosing the
polynomial in \eqref{eq:actualred} on a rectangular jet box remains valid
for the actual harmonic field contained in that box; the other box points
need not solve the harmonic equations.

The actual scalar Hessian is
\begin{equation}
 Q^\phi=2q^A{}_Dq^{DB}\nabla_A\nabla_B\phi_0
       -2q^{AB}\nabla_A\nabla_B\eta
       +2q^{AB}P^C{}_{AB}\partial_C\phi_0
       -\mathcal V'''\eta^2.                         \label{eq:scalar}
\end{equation}
The general expression additionally contains
$2q^C{}_D C_1^D\partial_C\phi_0-2C_1^C\partial_C\eta$.
Metric and scalar unit-principal scaled forcing are respectively
$-2\ell^2\rho^2Q^{E,\mathrm{red}}$ and $+\ell^2\rho^2Q^\phi$.
These signs include the retained reduced wave operator convention.

\subsection{Exact angular reconstruction}

Define $S=n_1^2+n_2^2$, $Z=S-1/2$ and $F=Y^2-S/2+1/12$.
The scalar basis is $(1,Z,F)$, the vector basis $(DZ,DF)$, and the angular
tensor basis is
\begin{equation}
 (\kappa,Z\kappa,F\kappa,\operatorname{Hess}^{\rm TF}Z,
 \operatorname{Hess}^{\rm TF}F,T_+,T_-).
\end{equation}
The two retained transverse-traceless orientations are independent basis
features, even when their particular source coefficients agree.
There are $3+3+2+3+2+7=20$ metric amplitudes and three scalar amplitudes.
The inherited homogeneous boundary operator retains five Dirichlet and
eighteen Robin rows; no row is replaced by a longitudinal ansatz here.

The exact exporter leaves all 42 ordinary metric two-jet entries and six
scalar two-jet entries independent. Rational rank witnesses extract
coefficients, after which complete polynomial reconstruction modulo
$\sum_i n_i^2-1$ is mandatory. Thus finite witnesses alone are not the proof.
Laurent coefficients are compared exactly, retaining both TT sectors and
the general gauge corrections. An ambient frame with four sphere labels and two normal/time labels is
used with tangential projection: the physical spacetime dimension is still
five. The machine-readable maps supplement these
formulas; printing thousands of monomials is unnecessary for replay.

\section{Actual first jets and source enclosures}

The predecessor supplies actual scaled $C^3$ master/gauge approximation
errors and background radial jet enclosures through order four~\cite{parent}.
The new oracle
composes them with explicit comparison functions, normalized bivariate
Taylor jets, frame factors and radial conversion. In particular,
\begin{equation}
 \partial_t=\ell^{-1}\partial_x,\quad
 \partial_y=-(\ell\rho)^{-1}\partial_r,\quad
 F_{yy}=(\ell\rho)^{-2}F_{rr}+a(\ell\rho)^{-1}F_r.
\end{equation}
The metric two-jet requires gauge amplitudes through order three; the scalar
two-jet requires the master through order three. These are actual derivative
error certificates, not derivatives inferred from a value-error bound.
The comparison functions are retained in shared expressions before interval
evaluation. The actual solution is not replaced by a comparison.

On backward cones wholly in $r\ge x$, the dynamic shell-lift term vanishes;
a separately checked resolvent recomputation removes only that contribution
from the exterior error budget. The full error budget is used elsewhere.
At $r=0$ the known physical boundary identities are imposed exactly.

All acceptance inequalities are rational with outward rounding at denominator
$2^{112}$. A source-only integer endpoint implementation uses the same
floor/ceiling interval operations and is checked against the Fraction
implementation. Its use changes replay cost, not the enclosures or the
mathematical premises. No floating-point plot controls a sign decision.

The domain cover consists of closed mesh boxes intersected with the triangular
domain in \eqref{eq:domain}. Evaluation on their outer rectangles only widens
explicit comparison and background expressions; inherited actual errors are
asserted only at points of the intersection. No extension of the PDE solution
outside $D$ is presumed. The radial and shell profiles also use complete cells.
Define the full bulk forcing and its uniform angular bundle bound by
\begin{equation}
 \mathcal F=(-2\ell^2\rho^2Q^{E,\mathrm{red}},\,+\ell^2\rho^2Q^\phi),
 \qquad F_H=\sup_{D_{\rm base}}|\mathcal F|_H.
\end{equation}
The enclosure table is generated and checked directly from the rational reports:
\begin{center}
\begin{tabular}{lll}
\toprule
Complete cell & Coefficient & Outward enclosure \\
\midrule
Exterior bulk & $\langle Q^E_{yy}\rangle$ & $[-6.12,-5.03]\,10^{17}$ \\
Reflected bulk & $\langle Q^\phi\rangle$ & $[-1.28,-1.03]\,10^{11}$ \\
Physical shell & $\langle Q_p\rangle$ & $[4.63,5.72]\,10^8$ \\
\bottomrule
\end{tabular}
\end{center}
Here $\delta=1/4096$. The exterior cell is $[1/4,1/4+\delta]\times[5/4,5/4+\delta]$ in $(x,r)$; the reflected cell is $[3/2,3/2+\delta]\times[1/4,1/4+\delta]$; the shell cell is $[5/4,5/4+\delta]\times\{0\}$. These are closed-cell enclosures in model units.

The full collar is covered by 81 clipped boxes, with a further 64-cell profile on $r\in[1,2]$, $x\in[1/4,1/4+\delta]$, and six complete shell time cells. Every profile cell has strictly negative natural and reduced $yy$ mean coefficients. The global bound is
\[F_H\le17\,387\,513\,156\,573<171\,610\,068\,644\,000,\]
where the rightmost number is the retained upper bound. Their ratio is greater than $9.8$. The rounded response majorant $1.03\,10^{33}$ and contraction $0.002109$ used below are retained.


The table's $Q_p$ is the one-copy outward-normal scalar junction Hessian.
Its two-sided jump and inward scaled forcing require the separate retained
copy/orientation factors; they are not the quantity printed in that row.

The source coefficients in that table are frame/basis components. Their signs
alone do not describe a physical force. The global bulk forcing bound uses
the positive metric Frobenius norm and probability-normalized angular $L^2$
norm. Improving that forcing bound does not itself update the shell Green
operator, contraction constant or a numerical response residual.

\section{An intrinsic shell scalar and certified mean zeros}

Let $X_\eps$ embed the material shell, $\gamma_\eps=X_\eps^*g_\eps$ and
$\psi_\eps=X_\eps^*\phi_\eps$. Define the Lorentzian intrinsic scalar
\begin{equation}
 I_\eps=\gamma_\eps^{\mu\nu}
             (d_\mu\psi_\eps)(d_\nu\psi_\eps).
\end{equation}
The background pullback is constant. Therefore $I_0=I'_0=0$ and
\begin{equation}
 I''_0=2\gamma_0^{\mu\nu}
             (d_\mu\eta_{\rm mat})(d_\nu\eta_{\rm mat}).
 \label{eq:intrinsic}
\end{equation}
No first/second metric variation or second scalar variation enters this order.
The field and embedding must be transformed jointly. In a normal displacement
description $\eta_{\rm mat}=\chi_b+p_b\zeta$; the plus-pullback changes
$\chi_b\mapsto\chi_b+p_b f$ and $\zeta\mapsto\zeta-f$ preserve it.
In the registered fixed-shell representation $\zeta=0$ and $f|_b=0$, so
$\eta_{\rm mat}=vY$. An additional physical embedding perturbation would
change this identification. The familiar second gauge transformation law
has only Lie derivatives of $I_0$ and $I'_0$~\cite{bmm2006,bmm2007}.
Hence \eqref{eq:intrinsic} is invariant under the joint perturbative
field/embedding transformations and shell relabellings in these premises.
The distinction between ambient and intrinsic hypersurface gauge follows
the standard geometric framework~\cite{mars}.

Since $|DY|^2=4(S-Y^2)$, the explicit coefficient and its probability mean are
\begin{align}
 I''_0&=\frac2{\rho_b^2}
       \left[-v_t^2Y^2+\frac{4v^2}{c^2}(S-Y^2)\right],\\
 M(t):=\avg{I''_0}&=\frac{-v_t^2+8v^2/c^2}{3\rho_b^2}.
 \label{eq:mean}
\end{align}
The exact moments are $\avg{Y^2}=1/6$ and $\avg{|DY|^2}=4/3$.
This average uses the chosen background time foliation; no physical clock
or observational sampling protocol is newly specified.

\begin{proposition}[Two simple mean zeros per certified bracket]
For the registered actual first solution, $M$ has exactly two simple zeros
in each of the scaled-time brackets $[5/4,11/8]$ and $[13/8,7/4]$ at $r=0$.
It is positive between the two zeros and negative on the remaining portions
of each bracket. In particular four distinct simple zeros are certified;
additional zeros outside these brackets are not excluded.
\end{proposition}
\begin{proof}
Where $v\ne0$ put $H=cv_t/v$ and
$N=vv_{tt}-v_t^2+\theta vv_t$. Direct differentiation gives
\begin{equation}
 H_t=\frac{cN}{v^2},\qquad
 M=\frac{v^2}{3\rho_b^2c^2}(8-H^2).
\end{equation}
The actual two-jet interval oracle proves $v<0$ throughout the first bracket,
$v>0$ throughout the second, and $N<0$ throughout both. On each bracket
$H$ is therefore strictly decreasing. Its left endpoint is greater than
$\sqrt8$ and its right endpoint less than $-\sqrt8$; these comparisons
are checked by signed squared rational inequalities. Continuity and strict
monotonicity give exactly one crossing of each level. Both crossings are
simple because $H_t<0$ and $H=\pm\sqrt8\ne0$ there. The positive prefactor
in the expression for $M$ proves the stated signs.
\end{proof}

At times with $vv_t\ne0$, the pointwise quadratic scalar itself has both
angular signs: it is negative at $(1,0,0,0)$ and positive at
$(1/\sqrt2,1/\sqrt2,0,0)$. For $S>0$ its angular zeros satisfy
$Y^2/S=4v^2/(c^2v_t^2+4v^2)$. A negative mean does not imply negativity
at every angle. These are statements about a Lorentzian squared gradient
at quadratic order, not attraction or repulsion. A finite-amplitude sign
would require an actual nonlinear family and control of its remainder.

\section{A quantitative evolved second-mean consequence}

The retained conditional solution $U=(k,w)\in C^1(D_{\rm base};H)$ obeys
\begin{equation}
 \|U\|_X=\max\left(\frac{\sup|U|_H}{3},
                         \sup|U_x|_H,\sup|U_r|_H\right)<1.03\,10^{33}.
\end{equation}
For $m(x,r)=\int_{S^3}w\,d\mu$, the mean functional has norm one since
$\mu(S^3)=1$ and $W_0=1$. Consequently
$|m_x|\le |U_x|_H<1.03\,10^{33}$; no angular point-evaluation factor is
needed and the value weight three does not enter this derivative bound.

\begin{corollary}[Persistence on two complete rectangles]
For $1/1000\le x\le1/200$, the actual selected formal coefficient satisfies
\begin{align}
 m(x,r)&>1.85\,10^{30},&73/64\le r\le585/512,\\
 m(x,r)&<-1.85\,10^{30},&951/512\le r\le119/64.
\end{align}
At every such time it has at least one radial mean zero in
$(585/512,951/512)$. The physical time interval is
$10^{-9}\le t\le5\,10^{-9}$ in the model's units.
\end{corollary}
\begin{proof}
The authenticated initial mean has strict magnitude greater than
$7\,10^{30}$ with the indicated signs over the two whole cells~\cite{parent}.
The fundamental theorem of calculus gives
$|m(x,r)-m(0,r)|<x\,1.03\,10^{33}\le5.15\,10^{30}$.
Subtracting this drift preserves the stated strict margin. Each rectangle
lies in $D$. Radial continuity supplies the zero separately at every time;
it does not prove uniqueness or a continuous zero branch.
\end{proof}

This is a quantitative corollary of an accepted derivative bound, not an
independent waveform calculation or a gauge-invariant observable. It retains
the bulk, shell and coupling effects through its global premise rather than
replacing the response with a local free wave.

\section{Verification and remaining work}

The capsule supplies the exact evaluator, projected coefficient maps,
interval oracle, complete-cell certificates, proofs, normal/optimized
goldens and reason-specific corrupted-formula controls. Independent internal
checks compare complete coordinate metric-family expansions with the
covariant frame calculation, reconstruct every angular polynomial, and
independently audit interval arithmetic, domain coverage, angular Gram
norms and the shell-zero inequalities. Complete reports are compared as
typed JSON and as raw bytes. The archive manifest and external reproduction
receipt state which programs are replayed and which older results remain
authenticated premises. Internal computational review is not external
specialist peer review or an assessment of originality.

The universal Ricci and gauge identities are already in the literature
\cite{bmm2006,bmm2007,xpert}; the vanishing-lower-order gauge criterion is
also established. This note makes no external novelty or priority claim.
Its concrete contribution is the completed model-specific source calculation,
actual interval enclosures, intrinsic diagnostic and quantitative corollaries.
The focused primary comparison is included with the source.

The next task is to evaluate a continuous selected-space numerical response
candidate and bound its full causal integral-equation residual, including
source/input, quadrature, reflection, initial, shell and reconstruction errors.
For a candidate $V$ and the retained contraction $q_*=0.002109$, a complete
residual $R=V-b-TV$ would give
$\|U-V\|_X\le\|R\|_X/(1-q_*)$.
Neither a Neumann tail count nor a source enclosure alone certifies that
residual. No first derivative of the quadratic source is asserted from the
available first-field $C^3$ bounds. No exact finite-amplitude nonlinear family,
observational fit or higher-dimensional origin of the Big Bang is established.

\begin{thebibliography}{9}
\bibitem{parent} R. Maldonado, \emph{Numerical response inputs and selected-space
evaluation}, reproducible conditional working checkpoint (2026),
\href{https://github.com/maldonado-research/HDblast/tree/main/research/HDBLAST_NUMERICAL_RESPONSE_INPUTS_20261016}{public numerical-input checkpoint}.
\bibitem{bmm2006} D. Brizuela, J. M. Mart\'in-Garc\'ia and G. A. Mena Marug\'an,
\emph{Second and higher-order perturbations of a spherical spacetime},
Phys. Rev. D \textbf{74}, 044039 (2006),
\href{https://arxiv.org/abs/gr-qc/0607025v1}{arXiv:gr-qc/0607025v1}.
\bibitem{bmm2007} D. Brizuela, J. M. Mart\'in-Garc\'ia and G. A. Mena Marug\'an,
\emph{High-order gauge-invariant perturbations of a spherical spacetime},
Phys. Rev. D \textbf{76}, 024004 (2007),
\href{https://arxiv.org/abs/gr-qc/0703069v1}{arXiv:gr-qc/0703069v1}.
\bibitem{xpert} D. Brizuela, J. M. Mart\'in-Garc\'ia and G. A. Mena Marug\'an,
\emph{xPert: Computer algebra for metric perturbation theory},
Gen. Relativ. Gravit. \textbf{41}, 2415--2431 (2009),
\href{https://arxiv.org/abs/0807.0824v1}{arXiv:0807.0824v1}.
\bibitem{mars} M. Mars, \emph{First and Second Order Perturbations of
Hypersurfaces}, Class. Quantum Grav. \textbf{22}, 3325--3348 (2005),
\href{https://arxiv.org/abs/gr-qc/0507005v1}{arXiv:gr-qc/0507005v1}.
\end{thebibliography}
\end{document}
